\documentclass[quick,con]{debate}
\motion{超级英雄的出现，对世界而言}
\stance{反方}{更是灾难}
\begin{document}
\debatetitle
\begin{thesisbox}[一句话立论]超级英雄的出现，第一次把一种压倒性的力量放在了没有任何人能够纠错的位置上，世界从此失去的不是安全感，是出错以后改回来的能力。\end{thesisbox}
\section{定义与判准（标准表述）}
\begin{bullets}
\item 超级英雄：拥有远超常人的能力，主动用这种能力介入公共安全事务；行动依据是自己的判断，不由现有法律和权力机构授权，也不由它们问责。三条缺一不可，第三条是全场的支点。
\item 出现：这种人存在于世界上成为事实，并且被世界知道。不是超能力普及到所有人。
\item 灾难：世界整体处境重大而且难以逆转地变糟，核心是纠错能力的丧失。不是“只要有风险”。
\item 判准：与正方同一版。长期看普通人的处境是变好还是变坏，看安全与自主两项。
\item 我方要证明：秩序代价是结构性的、改不回来的，且大到把救援抵消。对方要证明：他挡下的伤害真实重大，代价没有大到抵消它。
\item 对方提偏向定义或判准时的 30 秒口径：「对方想把这道题算成‘救了几个人’，那是半本账。判准里写着‘自主’这一项，纠错能力就在这一项里。请对方给这一栏记上代价，空着不填不叫赢，叫没算完。」
\end{bullets}
\section{我方论点}
\begin{longtable}{@{}L{\dimexpr0.050\linewidth-1.50\tabcolsep}L{\dimexpr0.156\linewidth-1.50\tabcolsep}L{\dimexpr0.499\linewidth-1.50\tabcolsep}L{\dimexpr0.295\linewidth-1.50\tabcolsep}@{}}
\toprule \thead{标签} & \thead{一句话} & \thead{核心数据 / 学理 / 案例} & \thead{出处}\\\midrule\endhead
没人能拦 & 压倒性的力量落在了没有任何现存机制能纠错的位置上 & 韦伯 1919：使用物理暴力的权利只在国家允许的限度内被赋予个人；印度护牛义警从2015年5月到2018年12月至少44人被打死、36 名穆斯林、约280人受伤、20 个邦100多起，近三分之一致死案件警方反过来追究受害方；凤凰侠2011年西雅图胡椒喷雾事件 & Weber《以政治为业》；HRW 2019；CBS News 2011，均已核实\\
坏人也会有 & 世界从此必须按最坏的持有者布防，这种不对称赔不起 & SIPRI 2025年鉴：2025年1月全球约12241枚核弹头，9614 枚可用、3912 枚已部署；1983 年9月26日苏联预警误报五枚导弹，起因是阳光在高云上的反射 & SIPRI Yearbook 2025 ch.6；Britannica “Stanislav Petrov”，均已核实\\
等靠退化 & 制度和普通人一起后退，等他不在时世界已经没有备份 & V-Dem 2025：45 国正在威权化，其中27国起点是民主国家，到2024年只剩9国仍是，死亡率67\%；全球72\%人口生活在威权国家；Nelson \& Norton 2005 研究三：112 人，报名率 23\% 对 42\%，90 天后到场 4\% 对 17\% & V-Dem《2025 年民主报告》；JESP 41(4):423–430，均已核实\\
\bottomrule\end{longtable}
\section{对方论点 → 一句话反驳}
\begin{longtable}{@{}L{\dimexpr0.265\linewidth-1.33\tabcolsep}L{\dimexpr0.312\linewidth-1.33\tabcolsep}L{\dimexpr0.423\linewidth-1.33\tabcolsep}@{}}
\toprule \thead{对方会说} & \thead{我方一句话回} & \thead{对方追问时}\\\midrule\endhead
救得下来，他补上了能力缺口 & 现实里的瓶颈从来不是能力，是授权和意愿 & 联合国自己说，受制于会员国不愿加强授权、不愿增派部队\\
恶做不成，他抬高了被抓住的确定性 & 你引的研究里全是可被投诉和起诉的警察 & 同样逻辑用在他身上，确定性就成了他一个人的心情\\
人会跟着做，榜样拉低门槛 & 同一篇论文的另一半直接写着相反的结果 & 90 天后到场 4\% 对 17\%，作者事先就是这样预测的\\
现实义警和他不一样，类比不成立 & 类比落点是能力与问责的落差，不是能力大小 & 没有超能力的人这个落差已经造成 44 条人命\\
\bottomrule\end{longtable}
\section{必问与必答}
\begin{fields}
\field{必问 （对方不答就点名）}{
\begin{enumerate}[leftmargin=1.8em,itemsep=2pt,topsep=2pt]
\item 他错了以后，谁能强制他停下来？（一辩稿抛出，全场追四次）
\item 判准里“自主”那一栏，你方今天记了多少代价？
\item 卢旺达缺的是能力还是授权？请念联合国那句话。
\item 你引用的威慑研究，研究对象可不可以被投诉和起诉？
\end{enumerate}
}
\field{必答}{
\begin{bullets}
\item “你说的灾难是和哪个世界比？” → 和他出现之前的世界比。之前的世界会失败，但失败能追责、能改制度。
\item “你是不是主张他一定会作恶？” → 不主张。我方的问题是无法验证，也无法在他错的时候改回来。
\item “那些本来会死的人怎么办？” → 我方承认他们的命是真的，也请对方承认另一栏的代价是真的。
\item “你引的研究没有对照组？” → 是，作者自己写明了，这一点我方先说的。主证据是制度那一组数据。
\end{bullets}
}
\end{fields}
\section{自由辩战场概要}
\begin{longtable}{@{}L{\dimexpr0.133\linewidth-1.50\tabcolsep}L{\dimexpr0.261\linewidth-1.50\tabcolsep}L{\dimexpr0.239\linewidth-1.50\tabcolsep}L{\dimexpr0.367\linewidth-1.50\tabcolsep}@{}}
\toprule \thead{战场} & \thead{开场问题} & \thead{目标} & \thead{收束语}\\\midrule\endhead
谁能把它改回来 & 他错了以后，谁能强制他停下来？ & 让评委听见正方第四次给不出办法 & 一整场比赛，对方没有给出一个能执行的办法\\
能力还是决心 & 卢旺达缺的是部队还是授权？ & 把正方最好的案例改判成授权缺口 & 对方最有力的两个案例，缺的都不是力气\\
外推是否成立 & 威慑研究的对象可不可以被起诉？ & 拆掉正方论据的适用范围 & 同一套逻辑放到他身上，确定性就成了他一个人的确定性\\
\bottomrule\end{longtable}
时间分配 90 秒 / 90 秒 / 60 秒，留 20 秒备用。站位：一辩守定义判准，二辩接驳论，三辩接盘问，四辩收束。

\section{如果……就……}
\begin{contingency}
\ifthen{正一没说“不受问责”，或对方抛出没预料到的数据}{前者由反四第一条链挖出来，反二驳论第一句点破。后者先问出处与年份，再问研究对象可不可以被问责。}
\ifthen{论点三被打穿（对照组问题被抓），或对方回避必问}{承认局限，收缩到制度侧数据（27 国只剩 9 国），自主项的损失仍然成立。回避就点名、重复、不换题，反三小结里说明这是第几次没有回答。}
\ifthen{全队红线}{不许说“他迟早会作恶”；不许把漫画反派当论据；引用 Nelson \& Norton 时必须同时说出“没有中性对照组”。}
\end{contingency}
\end{document}
